\documentclass[11pt,a4paper]{article}

% ── Codificación, fuentes y lenguaje ───────────────────────────────────────────
\usepackage[utf8]{inputenc}
\usepackage[T1]{fontenc}
\usepackage[default,scale=0.95]{sourcesanspro}
\renewcommand{\familydefault}{\sfdefault}
\usepackage[spanish,es-noshorthands,es-tabla]{babel}

% ── Geometría y espaciado ─────────────────────────────────────────────────────
\usepackage[top=2.2cm, bottom=2.5cm, left=2.2cm, right=2.2cm, headheight=15pt]{geometry}
\usepackage{setspace}
\setstretch{1.18}
\usepackage{parskip}

% ── Matemáticas y unidades ────────────────────────────────────────────────────
\usepackage{amsmath}
\usepackage{amssymb}
\usepackage{siunitx}

% ── Circuitos ─────────────────────────────────────────────────────────────────
\usepackage[european, straightvoltages]{circuitikz}

% ── Tablas ────────────────────────────────────────────────────────────────────
\usepackage{booktabs}
\usepackage{tabularx}
\usepackage{array}
\usepackage{longtable}
\usepackage{needspace}
\usepackage{colortbl}
\usepackage{enumitem}

% ── Colores, cajas e iconos ───────────────────────────────────────────────────
\usepackage[dvipsnames]{xcolor}
\usepackage{tcolorbox}
\tcbuselibrary{skins, breakable}
\usepackage{fontawesome5}
\usepackage{graphicx}

% ── Listados de código (dark-mode infográfico) ────────────────────────────────
%   Estilo base: fondo oscuro con números de línea. Los bloques lstlisting se
%   envuelven automáticamente en una caja tcolorbox redondeada.
\usepackage{listings}

% ── Secciones con barra de color (infografía) ────────────────────────────────
\usepackage{titlesec}
\titleformat{\section}
  {\normalfont\LARGE\bfseries\color{azulNoche}}
  {\colorbox{cyanNeon}{\color{fondoOscuro}\thesection}}{0.7em}{}
  [\vspace{2pt}\color{cyanNeon}\rule{\linewidth}{1.8pt}]
\titlespacing*{\section}{0pt}{24pt}{10pt}
\titleformat{\subsection}
  {\normalfont\large\bfseries\color{azulNoche}}
  {\textcolor{cyanNeon}{\thesubsection}}{0.7em}{}
  [\vspace{1pt}\color{grisLinea}\rule{\linewidth}{0.6pt}]
\titlespacing*{\subsection}{0pt}{14pt}{6pt}
\titleformat{\subsubsection}
  {\normalfont\normalsize\bfseries\color{azulNoche}}
  {\textcolor{cyanNeon}{\thesubsubsection}}{0.7em}{}
\titlespacing*{\subsubsection}{0pt}{10pt}{4pt}

% ── Cabeceras y pies de página ────────────────────────────────────────────────
\usepackage{fancyhdr}
\usepackage{lastpage}
\pagestyle{fancy}
\fancyhf{}
\renewcommand{\headrulewidth}{0pt}
\renewcommand{\footrulewidth}{0pt}
\fancyfoot[C]{%
  \begin{minipage}{\linewidth}
    \vspace{2pt}
    \color{cyanNeon}\rule{\linewidth}{1.2pt}\\[2pt]
    \centering\color{grisTexto}\footnotesize
    Página \thepage\ de \pageref{LastPage}
  \end{minipage}%
}

% ── Hiperenlaces ──────────────────────────────────────────────────────────────
\usepackage[colorlinks=true, linkcolor=azulNoche, urlcolor=cyanNeon]{hyperref}
\sloppy
\emergencystretch 3em

% ── Paleta de colores — tecnológico dark-mode ──────────────────────────────────
\definecolor{fondoOscuro}{HTML}{0D1B2A}%             Dark navy — fondo banda título
\definecolor{verdeTurquesa}{HTML}{004D40}%           Verde turquesa — bandas de marca
\definecolor{azulNoche}{HTML}{1B3A6B}%               Midnight blue — secciones, marcos
\definecolor{azulMedio}{HTML}{2A5298}%               Azul medio — variante de acento
\definecolor{cyanNeon}{HTML}{00C8E0}%                Cyan eléctrico — acentos primarios
\definecolor{verdeSignal}{HTML}{00B887}%             Verde señal — fortalezas, éxito
\definecolor{naranjaVivo}{HTML}{FF7043}%             Naranja vivo — mejoras, advertencias
\definecolor{rojoAlerta}{HTML}{E53935}%              Rojo alerta — errores
\definecolor{grisPapel}{HTML}{F4F6FA}%               Fondo tarjetas claras
\definecolor{grisLinea}{HTML}{C8D0E0}%               Bordes sutiles
\definecolor{grisTexto}{HTML}{6B7A99}%               Texto secundario
\definecolor{amarilloNota}{HTML}{FFB300}%            Notas / advertencias doradas
\definecolor{codigoFondo}{HTML}{1E2D3D}%             Fondo del bloque de código
\definecolor{codigoComentario}{HTML}{7A8FA6}%        Comentarios
\definecolor{codigoCadena}{HTML}{FF9D5C}%            Cadenas / strings
\definecolor{codigoKeyword}{HTML}{00C8E0}%           Palabras clave
\definecolor{codigoNumero}{HTML}{00E5A0}%            Números

% ── Banda de título: portada infográfica ──────────────────────────────────────
%   Diseño en 2 capas: banda de color (por defecto fondo oscuro) + franja de
%   acento cyan abajo. Uso: \bandaTitulo[color]{título}{subtítulo}.
%   El icono se puede cambiar con \renewcommand{\iconoBanda}{...} (FontAwesome).
\newcommand{\iconoBanda}{microchip}
\newcommand{\bandaTitulo}[3][fondoOscuro]{%
  \begin{tcolorbox}[
    enhanced,
    colback=#1, colframe=#1,
    boxrule=0pt, arc=8pt, outer arc=8pt,
    width=\linewidth,
    top=18pt, bottom=6pt, left=14pt, right=14pt,
    borderline south={3.5pt}{0pt}{cyanNeon},
  ]
    \begin{minipage}[c]{0.07\linewidth}
      \centering
      \textcolor{cyanNeon}{\fontsize{30}{30}\selectfont\faIcon{\iconoBanda}}
    \end{minipage}%
    \hspace{8pt}%
    \begin{minipage}[c]{0.88\linewidth}
      {\fontsize{20}{24}\selectfont\bfseries\color{white}#2}\par\vspace{5pt}%
      \textcolor{cyanNeon!80!white}{\small\faIcon{angle-right}~#3}%
    \end{minipage}
    \vspace{4pt}
  \end{tcolorbox}%
  \vspace{8pt}%
}

% ── Tarjetas de datos (infografía) ────────────────────────────────────────────
\tcbset{
  tarjetaDato/.style={
    enhanced, breakable,
    arc=6pt, outer arc=6pt,
    colback=grisPapel, colframe=azulNoche,
    boxrule=1pt,
    fonttitle=\bfseries\small, coltitle=white,
    attach boxed title to top left={yshift=-3mm, xshift=6mm},
    boxed title style={
      colback=azulNoche, colframe=azulNoche,
      arc=4pt, boxrule=0pt,
    },
    drop fuzzy shadow=azulNoche!30!white,
    top=8pt, bottom=8pt, left=10pt, right=10pt,
  },
  tarjetaIcono/.style={
    enhanced, breakable,
    arc=6pt, outer arc=6pt,
    colback=white, colframe=grisLinea, boxrule=0.8pt,
    fonttitle=\bfseries\small, coltitle=azulNoche,
    borderline west={3pt}{0pt}{cyanNeon},
    drop fuzzy shadow=azulNoche!25!white,
    top=6pt, bottom=6pt, left=10pt, right=10pt,
  },
  cajaTitulo/.style={
    enhanced, breakable,
    arc=6pt, outer arc=6pt,
    colback=azulNoche!10!grisPapel, colframe=azulNoche,
    boxrule=1pt,
    fonttitle=\bfseries, coltitle=white,
    attach boxed title to top left={yshift=-3mm, xshift=6mm},
    boxed title style={
      colback=azulNoche, colframe=azulNoche,
      arc=4pt, boxrule=0pt,
    },
    drop fuzzy shadow=azulNoche!25!white,
    top=8pt, bottom=8pt, left=10pt, right=10pt,
  },
  cajaContenido/.style={
    enhanced, breakable,
    arc=5pt, outer arc=5pt,
    colback=grisPapel, colframe=grisLinea,
    boxrule=0.8pt,
    fonttitle=\bfseries\small, coltitle=azulNoche,
    top=6pt, bottom=6pt, left=10pt, right=10pt,
  },
  cajaFortaleza/.style={
    enhanced, breakable,
    arc=6pt, outer arc=6pt,
    colback=verdeSignal!8!white, colframe=verdeSignal,
    boxrule=1pt,
    borderline west={4pt}{0pt}{verdeSignal},
    fonttitle=\bfseries\small, coltitle=verdeSignal!70!black,
    drop fuzzy shadow=verdeSignal!20!white,
    top=6pt, bottom=6pt, left=10pt, right=10pt,
  },
  cajaMejora/.style={
    enhanced, breakable,
    arc=6pt, outer arc=6pt,
    colback=naranjaVivo!8!white, colframe=naranjaVivo,
    boxrule=1pt,
    borderline west={4pt}{0pt}{naranjaVivo},
    fonttitle=\bfseries\small, coltitle=naranjaVivo!80!black,
    drop fuzzy shadow=naranjaVivo!20!white,
    top=6pt, bottom=6pt, left=10pt, right=10pt,
  },
  cajaRecomendacion/.style={
    enhanced, breakable,
    arc=6pt, outer arc=6pt,
    colback=cyanNeon!6!white, colframe=cyanNeon!60!azulNoche,
    boxrule=1pt,
    borderline west={4pt}{0pt}{cyanNeon},
    fonttitle=\bfseries\small, coltitle=azulNoche,
    drop fuzzy shadow=azulNoche!20!white,
    top=6pt, bottom=6pt, left=10pt, right=10pt,
  },
}

% ── Ícono + texto (encabezados de sección y elementos) ────────────────────────
\newcommand{\iconotexto}[2]{\textcolor{cyanNeon}{\faIcon{#1}}~#2}

% ── Listados de código: estilo dark + auto-envoltura ─────────────────────────
\lstdefinestyle{estiloCodigo}{
    backgroundcolor=\color{codigoFondo},
    basicstyle=\ttfamily\small\color{white},
    commentstyle=\color{codigoComentario}\itshape,
    stringstyle=\color{codigoCadena},
    keywordstyle=\color{codigoKeyword}\bfseries,
    numberstyle=\tiny\color{codigoComentario},
    numbers=left,
    numbersep=10pt,
    showstringspaces=false,
    breaklines=true,
    breakatwhitespace=false,
    tabsize=4,
    frame=none,
    captionpos=b,
    aboveskip=6pt,
    belowskip=4pt,
    xleftmargin=14pt,
    framexleftmargin=14pt,
    literate=
      {á}{{\'a}}1 {é}{{\'e}}1 {í}{{\'i}}1 {ó}{{\'o}}1 {ú}{{\'u}}1
      {Á}{{\'A}}1 {É}{{\'E}}1 {Í}{{\'I}}1 {Ó}{{\'O}}1 {Ú}{{\'U}}1
      {ñ}{{\~n}}1 {Ñ}{{\~N}}1 {ü}{{\"u}}1 {Ü}{{\"U}}1
      {¿}{{?`}}1 {¡}{{!`}}1
}
\lstdefinestyle{estiloPython}{
    style=estiloCodigo,
    language=Python,
}
\lstset{style=estiloCodigo}
\BeforeBeginEnvironment{lstlisting}{%
  \begin{tcolorbox}[
    enhanced, breakable,
    arc=6pt, outer arc=6pt,
    colback=codigoFondo, colframe=azulNoche,
    boxrule=1pt,
    drop fuzzy shadow=azulNoche!25!white,
    top=2pt, bottom=2pt, left=0pt, right=0pt,
  ]%
}
\AfterEndEnvironment{lstlisting}{\end{tcolorbox}}

% ── Cajas de contenido temáticas ──────────────────────────────────────────────
\newtcolorbox{cajaRecuerda}{
    enhanced, breakable,
    arc=6pt, outer arc=6pt,
    colback=azulNoche!10!grisPapel,
    colframe=azulNoche, boxrule=1pt,
    borderline west={4pt}{0pt}{cyanNeon},
    fonttitle=\bfseries\small\color{white},
    title={\faIcon{info-circle}~Recuerda},
    attach boxed title to top left={yshift=-3mm, xshift=6mm},
    boxed title style={colback=azulNoche, arc=4pt, boxrule=0pt},
    drop fuzzy shadow=azulNoche!25!white,
    left=10pt, right=10pt, top=8pt, bottom=8pt,
}

\newtcolorbox{cajaConcepto}{
    enhanced, breakable,
    arc=6pt, outer arc=6pt,
    colback=verdeSignal!8!white,
    colframe=verdeSignal, boxrule=1pt,
    borderline west={4pt}{0pt}{verdeSignal},
    fonttitle=\bfseries\small\color{white},
    title={\faIcon{lightbulb}~Concepto clave},
    attach boxed title to top left={yshift=-3mm, xshift=6mm},
    boxed title style={colback=verdeSignal!80!black, arc=4pt, boxrule=0pt},
    drop fuzzy shadow=verdeSignal!20!white,
    left=10pt, right=10pt, top=8pt, bottom=8pt,
}

\newtcolorbox{cajaEjemplo}{
    enhanced, breakable,
    arc=6pt, outer arc=6pt,
    colback=cyanNeon!5!white,
    colframe=cyanNeon!70!azulNoche, boxrule=1pt,
    borderline west={4pt}{0pt}{cyanNeon},
    fonttitle=\bfseries\small\color{fondoOscuro},
    title={\faIcon{code}~Ejemplo},
    attach boxed title to top left={yshift=-3mm, xshift=6mm},
    boxed title style={colback=cyanNeon, arc=4pt, boxrule=0pt},
    drop fuzzy shadow=azulNoche!20!white,
    left=10pt, right=10pt, top=8pt, bottom=8pt,
}

\newtcolorbox{cajaEjercicio}{
    enhanced, breakable,
    arc=6pt, outer arc=6pt,
    colback=azulMedio!7!white,
    colframe=azulMedio, boxrule=1pt,
    borderline west={4pt}{0pt}{azulMedio},
    fonttitle=\bfseries\small\color{white},
    title={\faIcon{pencil-alt}~Ejercicio},
    attach boxed title to top left={yshift=-3mm, xshift=6mm},
    boxed title style={colback=azulMedio, arc=4pt, boxrule=0pt},
    drop fuzzy shadow=azulNoche!20!white,
    left=10pt, right=10pt, top=8pt, bottom=8pt,
}

\newtcolorbox{cajaReto}{
    enhanced, breakable,
    arc=6pt, outer arc=6pt,
    colback=naranjaVivo!6!white,
    colframe=naranjaVivo, boxrule=1pt,
    borderline west={4pt}{0pt}{naranjaVivo},
    fonttitle=\bfseries\small\color{white},
    title={\faIcon{fire}~Reto},
    attach boxed title to top left={yshift=-3mm, xshift=6mm},
    boxed title style={colback=naranjaVivo, arc=4pt, boxrule=0pt},
    drop fuzzy shadow=naranjaVivo!20!white,
    left=10pt, right=10pt, top=8pt, bottom=8pt,
}

% ── Indicadores de dificultad ─────────────────────────────────────────────────
\newcommand{\facil}{\textcolor{verdeSignal}{\faIcon{circle}\,\textbf{Fácil}}}
\newcommand{\intermedio}{\textcolor{naranjaVivo}{\faIcon{circle}\,\textbf{Intermedio}}}
\newcommand{\dificil}{\textcolor{rojoAlerta}{\faIcon{circle}\,\textbf{Difícil}}}

% ─────────────────────────────────────────────────────────────────────────────
\renewcommand{\iconoBanda}{shield-alt}
\fancyhead[L]{\footnotesize\color{grisTexto}\faIcon{shield-alt}~CYBERSEGURIDAD}
\fancyhead[R]{\footnotesize\color{grisTexto}Semana 03}
\hypersetup{pdftitle={Ciberseguridad Aplicada},pdfauthor={Prof. César Rodríguez}
  pdfsubject={CYBERSEGURIDAD}}

\begin{document}
\bandaTitulo[fondoOscuro]{Semana 03 --- Tipificacion, admisibilidad y codigo de etica}{Grupo 4 $\cdot$ rol Legal y reporte $\cdot$ ejemplo de trabajo}

\begin{center}\small\color{grisTexto}Prof. César Rodríguez \textbar{} ELECTRONICA \textbar{} 1 de noviembre de 2026\end{center}

\begin{tabularx}{\linewidth}{@{}>{\raggedright\arraybackslash}p{4.6cm}X@{}}
\toprule
\textbf{Curso} & CYBERSEGURIDAD --- Ciberseguridad Aplicada \\
\textbf{Semana} & 03 de 10 \\
\textbf{Fase} & 1 --- Fundamentos, criptografía y marco legal \\
\textbf{Grupo} & 4 \\
\textbf{Rol} & Legal y reporte \\
\textbf{Entrega} & GitHub, carpeta semana\_03/ \\
\textbf{Corte} & Domingo 23:59 de la semana correspondiente; lo posterior no se corrige. \\
\bottomrule
\end{tabularx}

\needspace{5\baselineskip}
\section{\iconotexto{user-shield}{Tu rol en la práctica}}

Verifica que el alcance y la autorización cubran lo ejecutado, anonimiza la evidencia, revisa la cadena de custodia e integra el informe técnico y el resumen ejecutivo.

\begin{itemize}[leftmargin=1.4em,itemsep=2pt,topsep=3pt]
  \item Revisar que el alcance y la autorización cubran exactamente lo ejecutado, ni más ni menos.
  \item Anonimizar la evidencia y verificar la cadena de custodia de cada archivo entregado.
  \item Integrar informe técnico y resumen ejecutivo, y declarar los límites del trabajo.
\end{itemize}

\needspace{5\baselineskip}
\section{\iconotexto{flag}{Qué aporta este rol al entregable}}

La versión final del informe y la declaración explícita de lo que no se pudo verificar.

\needspace{5\baselineskip}
\section{\iconotexto{list-ol}{Ejemplo paso a paso}}

El rol Legal traduce el caso ficticio a obligaciones concretas: que conducta encaja en que tipo, que evidencia se sostiene y cual se contamino. La etica se firma de forma individual, no en nombre del grupo.

\subsection{Separar hechos, pruebas y conclusiones del caso}

Primero se transcribe lo que ocurrio, sin calificar. Despues se marca que parte tiene respaldo documental y que parte es inferencia. Mezclarlas es el error que hace inadmisible un informe.

\subsection{Tipificar sin inventar articulos}

El tipo se identifica por su denominacion y se cita el articulo solo despues de verificarlo en el texto de la norma. Si no se pudo verificar, se escribe "pendiente de verificar" y se explica que falta.

\begin{lstlisting}[style=estiloCodigo]
mkdir -p semana_03/evidencia
sha256sum semana_03/norma.pdf > semana_03/evidencia/norma.sha256
printf 'norma verificada: %s\n' "$(date -u +%FT%TZ)" \
  > semana_03/evidencia/verificacion_norma.txt
\end{lstlisting}

\subsection{Decidir que evidencia es admisible}

Se evalua cada pieza con tres preguntas: se obtuvo con autorizacion, se conservo su integridad y se puede rastrear su origen. Si falla una, se excluye y se dice por que.

\subsection{Firmar el codigo de etica de forma individual}

Cada estudiante firma su propia copia. La firma del grupo no sustituye la individual: la responsabilidad no se delega.

\begin{lstlisting}[style=estiloCodigo]
sha256sum semana_03/codigo_etica_firmado_*.pdf \
  > semana_03/evidencia/etica.sha256
\end{lstlisting}

\needspace{8\baselineskip}\subsection{Conducta, tipo aplicable y evidencia}

\begin{center}
\begin{minipage}{\linewidth}
\centering
{\footnotesize
\textbf{Conducta, tipo aplicable y evidencia}\par\smallskip
\begin{tabular}{@{}>{\raggedright\arraybackslash}p{3.20cm}>{\raggedright\arraybackslash}p{3.00cm}>{\raggedright\arraybackslash}p{2.60cm}>{\raggedright\arraybackslash}p{1.60cm}>{\raggedright\arraybackslash}p{3.60cm}@{}}
\toprule
\textbf{Conducta} & \textbf{Tipo aplicable} & \textbf{Evidencia} & \textbf{Admisible} & \textbf{Motivo} \\
\midrule
Acceso sin autorizacion & Acceso indebido a sistema protegido & Registro del servidor & No & Se obtuvo fuera del alcance \\
Copia de datos de clientes & Violacion de privacidad de datos & Archivo sin hash & No & Integridad no verificable \\
Prueba dentro del alcance & Actividad autorizada & Bitacora y hashes & Si & Autorizacion previa y trazabilidad \\
Captura publicada sin anonimizar & Infraccion del codigo de etica & Publicacion en repositorio & No & Expone datos de terceros \\
\bottomrule
\end{tabular}
}
\par\smallskip{\footnotesize El tipo se escribe por su denominacion; el numero de articulo solo si se verifico contra la norma. La ultima columna justifica la decision, no la repite.}
\end{minipage}
\end{center}

\needspace{5\baselineskip}
\section{\iconotexto{box-open}{Qué guardar en el repositorio}}

Se entregan el analisis del caso, la tabla de admisibilidad con motivo por pieza, la verificacion de la norma y el codigo de etica firmado por cada estudiante.

\begin{lstlisting}[style=estiloCodigo]
semana_03/
|-- autorizacion.md
|-- analisis_caso.md
|-- codigo_etica_firmado_1.pdf
|-- codigo_etica_firmado_2.pdf
|-- evidencia/
|   |-- norma.sha256
|   |-- verificacion_norma.txt
|   `-- etica.sha256
`-- README.md
\end{lstlisting}

\begin{itemize}[leftmargin=1.4em,itemsep=2pt,topsep=3pt]
  \item Cada tipo aplicable cita la fuente verificada o declara que falta verificarla.
  \item Cada pieza de evidencia tiene decision y motivo.
  \item Hay un codigo de etica firmado por cada integrante.
  \item El analisis distingue hecho, inferencia y conclusion.
\end{itemize}

\needspace{5\baselineskip}
\section{\iconotexto{ruler-combined}{Cómo se corrige}}

\begin{tabularx}{\linewidth}{@{}X r@{}}
\toprule
Criterio & Peso\\
\midrule
Alcance explícito con inclusiones y exclusiones & 3,0 \\
Reglas de enfrentamiento operativas & 2,5 \\
Tipificación correcta de las conductas del caso & 2,0 \\
Juicio fundamentado sobre admisibilidad de la evidencia & 1,5 \\
Código de ética firmado individualmente & 1,0 \\
\midrule
\textbf{Total} & \textbf{10,0}\\
\bottomrule
\end{tabularx}

\begin{itemize}[leftmargin=1.4em,itemsep=2pt,topsep=3pt]
  \item Publicar capturas con datos de terceros es una infraccion del codigo, no un detalle de estilo.
  \item La verificacion de la norma se hace sobre el PDF entregado, no sobre una fuente de memoria.
\end{itemize}

\begin{cajaRecuerda}El trabajo es independiente: cada grupo usa su propio repositorio y su propio laboratorio, y la comparación de resultados es en clase. Corte: Domingo 23:59 de la semana correspondiente; lo posterior no se corrige.
\end{cajaRecuerda}
\end{document}
